% SLAJD 4
\section{Streszczenie} 
\begin{frame}{Streszczenie i cel pracy}
    \begin{itemize}
        \item Tekst analizuje rosnący trend migracji z miast na wieś.
        \item Argumentuje się, że życie poza aglomeracją oferuje wyższą jakość codziennego funkcjonowania.
        \item Wskazane są cztery najważniejsze obszary przewagi wsi: 
        \begin{itemize}
            \item zdrowie,
            \item relacje społeczne,
            \item ekonomia,
            \item psychologia.
        \end{itemize}
        \item \textbf{Wniosek:} Ucieczka z miasta jest świadomym wyborem na rzecz odzyskania równowagi życiowej i polepszenia psychiki.
    \end{itemize}
\end{frame}

% SLAJD 5
\section{Wprowadzenie}
\begin{frame}{Miasto a wieś - obecne trendy}
    \begin{itemize}
        \item Miasto jest z natury bardziej atrakcyjne (bliskość pracy, edukacji, rozrywki).
        \item Wsie i mniejsze miejscowości zyskują jednak coraz więcej mieszkańców.
        \item Główny czynnik zmiany: możliwość pracy zdalnej.
    \end{itemize}
\end{frame}

% SLAJD 6
\begin{frame}{Przyczyny migracji}
    \begin{itemize}
        \item Wzrastające ceny mieszkań w miastach.
        \item Ucieczka od miejskiego hałasu.
        \item Chęć życia w spokoju i ciszy.
        \item Zniechęcające cechy miasta: ciągła pogoń, stanie w korkach, wysokie koszty życia.
    \end{itemize}
\end{frame}

% SLAJD 7
\begin{frame}{Wizualizacja}
    \begin{figure}
        \centering
        \includegraphics[height=0.7\textheight]{wies.png} 
        \caption{Sielski krajobraz wiejski}
    \end{figure}
\end{frame}

% SLAJD 8
\begin{frame}{Kontakt z naturą i zdrowie}
    \begin{itemize}
        \item Na wsi występuje znacznie bliższy kontakt z naturą (lasy, pola zamiast betonu).
        \item Czyste powietrze zamiast spalin.
        \item Korzyści zdrowotne:
        \begin{itemize}
            \item mniej alergii,
            \item niższy poziom stresu.
        \end{itemize}
        \item Naturalna motywacja do aktywności (spacery, rower) bez konieczności szukania zatłoczonego parku.
    \end{itemize}
\end{frame}

% SLAJD 9
\begin{frame}{Relacje społeczne}
    \begin{itemize}
        \item \textbf{Miasto:} Mimo tłumów, łatwo o anonimowość i samotność.
        \item \textbf{Wieś:} Społeczności są mniejsze i bardziej zżyte.
        \item Sąsiedzi się znają i pomagają sobie.
        \item Tworzenie poczucia wspólnoty daje większe bezpieczeństwo i oparcie w trudnych chwilach.
    \end{itemize}
\end{frame}

% SLAJD 10
\begin{frame}{Przestrzeń życiowa}
    \begin{itemize}
        \item Porównanie ekonomiczne: Małe mieszkanie z widokiem na blok (miasto) vs Dom z własnym ogrodem (wieś).
        \item Nowa jakość życia na wsi:
        \begin{itemize}
            \item miejsce na grilla,
            \item przestrzeń do zabawy dla dzieci,
            \item możliwość wypicia kawy na własnym tarasie w otoczeniu zieleni (brak miejskiego zgiełku).
        \end{itemize}
    \end{itemize}
\end{frame}